\documentclass[12pt]{article}

\usepackage[utf8]{inputenc}
\usepackage[french]{babel}
\usepackage{hyperref}

\title{Gestion de Projet Informatique\\
        TD 7 : Mise en situation}
\author{Aida HASHEMI BENI, Louisa SOUICI, Tasnime EL OUESLATI, Cérine BEN 
AHMED}
\date{Année 2025/2026}

\begin{document}

\maketitle
\newpage
\tableofcontents
\newpage

\section{Organisation}

\subsection{[A1] Constitution d'un groupe}

[À compléter]

\subsection{[A2] Dépôt Git en ligne}

\url{https://github.com/cezs787/Gestion-de-Projet-Informatique}

[À compléter]

\subsection{[A3] Intégration des contributions}

[À compléter]



\section{Mise en situation Git}

\subsection{[B1] Gestion par branches}

[À compléter]

\subsection{[B2] Intégration de version et tag}

[À compléter]

\subsection{[B3] Gestion des propositions de contribution}

[À compléter]



\section{Outils avancés}

\subsection{[B4] Intégration du site web}

[À compléter]

\subsection{[C1] Répartition des tâches et milestones}

[À compléter]

\subsection{[C2] Intégration continue}

[À compléter]



\section{Gestion avancée et conclusion}

\subsection{[C3] Gestion de submodule}

Nous avons intégré à notre projet un submodule Git afin de gérer une dépendance externe. 
Pour démontrer l'utilisation des commandes d'intégration, le \texttt{Projet-JS} fera office de 
sous-module dans notre dépôt principal. \\
L'usage de la fonctionnalité submodule peut se montrer pertinente, lorsque le projet en cours dépend 
du contenu d'un autre projets ; son accès et sa visualisation sont alors facilités.

\subsubsection*{Ajout du submodule}

Pour ajouter le submodule au repertoire principal, il faut avoir recours à la commande suivante :

\begin{verbatim}
git submodule add https://github.com/cezs787/Projet-JS extern/Projet-JS
\end{verbatim}

Laquelle a créé automatiquement un fichier \texttt{.gitmodules} contenant la configuration du 
submodule.


\begin{figure}[h!]
    \centering
    \includegraphics[width=.85\linewidth]{C3.png}
    \caption{Fichier \texttt{.gitmodules} après ajout du sous-module.}
\end{figure}

\newpage

\subsubsection*{Utilisation par un nouvel utilisateur}

Pour pouvoir acceder au submodule, il est nécessaire pour le nouvel utilisateur de cloner 
correctement le projet avec le submodule, il doit alors avoir recours à la commande suivante :

\begin{verbatim}
git clone --recursive https://github.com/<nom-du-depot>.git
\end{verbatim}

Si le dépôt a déjà été cloné sans l'option \texttt{--recursive}, l'initialisation se fait via :

\begin{verbatim}
git submodule update --init --recursive
\end{verbatim}

\subsubsection*{Mise à jour du submodule}

Pour synchroniser le submodule avec la dernière version disponible dans son dépôt d'origine :

\begin{verbatim}
git submodule update --remote
\end{verbatim}

\begin{figure}[h!]
    \centering
    \includegraphics[width=.95\linewidth]{C3_2.png}
    \caption{Ajout du submodule via la commande \texttt{git submodule add}.}
\end{figure}

\subsection{[C4] Gestion par l'outil Project GitHub}

[À compléter]

\subsection{Conclusion}

[À compléter]

\end{document}
